\section{Choosing the constants and completing the reduction}
\label{sec:parameters}

The construction and its analysis have identified all the required
properties of the fixed parameters. We now choose them in an order
that makes the graph reduction possible. The key point is that the
list bound $d$ is fixed before the layer count and the source soundness,
so it cannot depend on the Label Cover alphabets.

Fix a rational $0<\delta<1/3$ and a rational
$0<\eps<\delta/8$. Set $L=64$ and take $\beta,K$ from
Lemma~\ref{lem:influences} for $L,\eps$.
Choose a positive integer $m$ with $L/m<\beta/2$, then a rational
$0<\lambda<1$ with $\lambda K<\eps$.
The coefficient law \eqref{eq:coefficient-law} is now fixed.
Since its probability $p_m$ of a full coefficient is positive, choose
an integer $s\ge2$ with $(1-p_m)^s<\eps$.

Next choose $\gamma>0$ sufficiently small that
$s(2\gamma)^2/\beta^2<\eps$, and then choose an even positive integer
$P$ sufficiently large that, with $D=mP$,
\[
 \frac{2L}{D}<\gamma,
 \qquad \frac{2Ls}{P}<\eps.
\]
This establishes every inequality in \eqref{eq:parameter-bounds}.
Lemma~\ref{lem:subchain-lists} gives the alphabet-independent bound
\[
 d=\max\{1,J(sL,\gamma)(m+1)^s\}.
\]
Apply Lemma~\ref{lem:alignment} to $s,d,\eta=\eps$, obtaining
$r\ge s$ and a rational law $\mu$ on $\binom{[r]}s$.
Finally choose a rational $\sigma\in(0,1)$ with
\begin{equation}\label{eq:source-soundness}
 4^r d^2\sigma<\eps.
\end{equation}
Only now invoke Theorem~\ref{thm:label-cover} for this fixed $\sigma$.
Its constant alphabets may be large, but no preceding choice depends
on their sizes or on the input formula.

All objects used to run the reduction are finite integers or finite
rational probability tables. For each fixed $\delta$, choose them
once and incorporate their finite descriptions into one algorithm.
Neither a uniform runtime as $\delta\to0$ nor an algorithm computing
all constants from $\delta$ is asserted or needed. There is no
input-dependent advice. The real tolerances, continuous functions,
junta approximants, and alignments enter only the analysis.

\begin{proof}[Proof of Theorem~\ref{thm:main}]
It suffices to treat rational $0<\delta<1/3$: for a fixed real
threshold, choose a fixed rational $0<\delta_0<\delta$ and use the
reduction for $\delta_0$.
Fix rational $\delta$ and choose the constants as above, with
$8\eps<\delta$. Use the preprocessing and graph construction of
Section~\ref{sec:construction}.

The two preprocessing branches already have the required promises.
A formula with no remaining clauses is satisfiable and produces a
one-vertex graph. An empty clause certifies unsatisfiability and
produces $K_q$, where $q>1/\delta$ and hence $1/q<\delta$.
For every other input, Proposition~\ref{prop:encoding} gives a
deterministic polynomial-time construction of a nonempty finite simple
unweighted graph. If the formula is satisfiable, its Label Cover
instance has value one, and Lemma~\ref{lem:completeness} supplies a
proper three-coloring of every output vertex.

Suppose instead that the formula is unsatisfiable, so its Label Cover
value is at most $\sigma$. If the metric clique branch occurs, the
output again has independent-set density $1/q<\delta$.
Otherwise, let $\mathcal A$ be any nonempty independent set of the
output graph. Lemma~\ref{lem:good-pairs} and
\eqref{eq:source-soundness} bound the probability of an unaligned
chain--$B$ pair by $2\eps$.
At every aligned pair, Proposition~\ref{prop:aligned-bound} bounds
the conditional probability of sampling $\mathcal A$ by $6\eps$;
at other pairs it is at most one. The multiplicity expansion in
Proposition~\ref{prop:encoding} makes this sampling law exactly the
uniform law on output vertices. Hence
\[
 \frac{|\mathcal A|}{|V(G)|}
 \le 6\eps+2\eps
 =8\eps<\delta.
\]
The empty independent set satisfies the same strict bound because
the graph is nonempty. Maximizing over all independent sets proves
the required soundness promise.
\end{proof}
